\looseness=-1 While the algebraic properties of the hyperbolic Kac--Moody groups and algebras attracted some attention, there have been only a few mathematical results concerning their geometry, for example, the study of homogeneous or symmetric spaces associated to them~\cite{FreynHartnickHornKoehl17}.
These classes of objects are very well understood for finite-dimensional Lie groups and for affine Kac--Moody groups. It is hoped that the understanding of this geometry, called Kac--Moody geometry, will shed new light on algebraic and structural properties of these groups. As a~first step towards the goal of understanding the geometry of hyperbolic Kac--Moody algebras and groups, we show in this work that for
hyperbolic Kac--Moody groups~$G$ over~$\C$, the associated Tits building is not only an abstract simplicial complex admitting an action of~$G$, but that it admits a natural embedding inside the compact real form~$\k$ of the Kac--Moody algebra~$\fg$. As a consequence, the structure of the Tits building and of the compact real form of the Kac--Moody algebra are closely related.

Our results generalize work of Quillen and Mitchell in the finite-dimensional case, and Kramer and Freyn in the affine case. Mitchell, in a paper based on ideas of Quillen, used embeddings of spherical buildings associated to simple (real or complex) Lie groups into the tangent space of
the associated noncompact real symmetric space~\cite{Mitchell88}. In particular, the topological building of a real noncompact Lie group~$G$ with maximal compact subgroup~$K$ is canonically identified with a space homeomorphic to the unit sphere in the tangent space of the noncompact real symmetric space~$G/K$. For example, the topological building of type~$A_1$, isomorphic to~$S^1$, can be embedded into the unit circle in the tangent space of $\mathcal{H}^2={\rm SL}_2(\R)/{\rm SO}(2)$ which is $\R^2$.

Kramer gave a topological construction of the complex twin building of type $A_n^{(1)}$ and an equivariant embedding of this building into the associated affine Kac--Moody algebra $\g$~\cite{Kramer02}.
Freyn gave a 2-parameter family $\varphi_{\ell,r}$ of equivariant embeddings of affine `twin cities' \mbox{$\cB \!=\! \cB^+\!\cup\! \cB^-$} into the `$s$-representations' of affine Kac--Moody symmetric spaces~\cite{Freyn09, Freyn10b, Freyn10d, Freyn14, Freyn15b}. A `twin city' is the natural completion of a twin building, so affine twin cities correspond to completions of affine Kac--Moody groups as twin buildings correspond to minimal Kac--Moody groups. Twin cities carry a natural topology that is derived from the topology on the corresponding Kac--Moody group.

Denoting by $\g=\k \oplus \p$ the Cartan decomposition, and restricting Freyn's result to non-completed affine Kac--Moody algebras, yields an identification of the twin building with the intersection $\p_{\ell, r}$ of the sphere of squared length $\ell\in \R$ with horospheres parametrized by $r_d=\pm r\not=0$, where $r_d$ is the real coefficient of the derivation~$d$. The positive and negative components of the twin building, $\cB^+$ and $\cB^-$, are immersed into the two sheets of $\p_{\ell, r}$ described by $r_d<0$ respectively~$r_d>0$.

For affine Kac--Moody groups of the compact type, which are symmetric spaces called of \emph{`type~II'} following Helgason's classification (see~\cite{Freyn15, Helgason01}), this result includes embeddings of complex buildings into the compact forms of affine Kac--Moody algebras; this case is the affine counterpart to the embeddings constructed in this paper for symmetrizable hyperbolic Kac--Moody algebras. We refer the reader to~\cite{Freyn09, Freyn10b, Freyn10a, Freyn10d} for additional details.

For $\g$ a hyperbolic Kac--Moody algebra we are motivated in part by the appearance of $\fg/\fk$ and $G/K$ in coset models of certain supergravity theories, where coset spaces of the split real forms occur as parameter spaces for the scalar fields of the theory \cite{DamourHenneauxNicolai02,Julia85, West01}.

The ideas in this paper may be extended to the wider class of Lorentzian Kac--Moody algebras, but it remains to be seen how much the results are affected by the differences in the geometry of the lightcone and the Tits cone. Some work in this direction was started by A.~Tichai~\cite{Tich14}.

\subsection{Summary of results}\label{subsection:Summary of results}

The following paragraph summarizes the main results of this paper.
For $G$ of hyperbolic type, we construct a $K$-equivariant embedding of the twin building of $G$ into the lightcone of the compact
real form $\fk$. In $\ft$, $W$ acts on rays in the interior of the lightcone (and on certain rays on the nullcone), tessellating a copy of
hyperbolic space $H^{n-1}$ in each half of the lightcone (and some limit points in its boundary), giving a twin apartment. The $W$-images of the fundamental domain form the chambers of that apartment, so each face of any chamber is in a hyperplane fixed by a $W$-conjugate of a simple reflection~$w_i$. The associated conjugate of the compact subgroup ${\rm SU}(2)_i$ fixes that hyperplane pointwise, but rotates the rest of $\ft$ into a $K$-conjugate of~$\ft$, another apartment sharing that fixed chamber wall.
This gives our geometric model of part of the twin building where each half is infinitely many copies of a tessellated hyperbolic space glued together along hyperplanes of the faces. Locally, at each such face, a family of chambers meet, the orbit of an ${\rm SU}(2)$ with a ${\rm U}(1)$ stabilizer, so the family is indexed by a Riemann sphere, ${\rm SU}(2)/{\rm U}(1) \cong S^2$.
